\section{Evidência de IMP-NOTIF-003 (``Limpar tudo'' por corte paginado)}

\textbf{Feature:} \texttt{IMP-NOTIF-003} --- marcar como lidas, em rodadas
paginadas e reentrantes, todas as notificações ativas do próprio UID.

\textbf{Estado anterior:} \texttt{NÃO IMPLEMENTADO}. Não havia endpoint; a única
mutação era a marcação individual de \texttt{IMP-NOTIF-002}.

\textbf{Estado depois:} \texttt{NÃO DIVERGENTE} para o escopo backend examinado.
O disparo pela interface pertence à caixa unificada (\texttt{IMP-NOTIF-001} e
\texttt{IMP-UI-004}), ainda abertas.

\subsection{Contrato}

Fontes: Seção~7 RN-M13-02; Seção~9 ``Marcar lida e Limpar tudo''; CUE
\texttt{\#M13LimparTudo}; Alloy \texttt{limparTudo},
\texttt{LimparTudoCorteEstavel}, \texttt{LoteSoProprias}, \texttt{SemDelete},
\texttt{WitnessLimparTudoRetry}, \texttt{WitnessLimparTudoCorteComNovaEmissao}.
Propriedades derivadas:

\begin{itemize}
  \item server-owned com autoridade M9 persistida relida na transação;
  \item marca apenas notificações endereçadas ao próprio UID
        (\texttt{id\_destinatario == uid}, RN-M13-01);
  \item \texttt{lida=true} com \texttt{lida\_em} do servidor; \texttt{lida=false}
        $\Rightarrow$ \texttt{lida\_em=null} (\#M13Notificacao);
  \item sem \texttt{DELETE}: o documento permanece para o histórico;
  \item corte estável: avisos emitidos após o corte ficam para a próxima
        invocação;
  \item retry/reentrância: avisos já marcados saem da consulta
        (\texttt{lida == false}) e não recebem novo instante.
\end{itemize}

A identidade M7 não se aplica como receipt de lote: \texttt{\#M13LimparTudo} e
\texttt{\#M13Marcacao} não carregam \texttt{id\_operacao}; a idempotência é
por item, preservando \texttt{lida\_em}, e a retomada é definida pelo corte
estável, não por uma operação única.

\subsection{Mudança}

\begin{itemize}
  \item \path{functions/src/notificacoes.ts}: nova callable
  \texttt{limparTudoNotificacoes}, com \texttt{corte} opcional (ISO 8601) e
  \texttt{limite} de página (padrão 100, máximo 200). A consulta filtra
  \texttt{id\_destinatario == uid}, \texttt{lida == false} e
  \texttt{emitida\_em <= corte}, ordena por \texttt{emitida\_em} e marca a
  página inteira na mesma transação; estado \texttt{lida}/\texttt{lida\_em}
  incoerente falha fechado;
  \item \path{firestore.indexes.json}: índice composto
  \texttt{Notificacoes(id\_destinatario, lida, emitida\_em)} referenciado por
  \path{firebase.json}.
\end{itemize}

\subsection{TDD}

\small
\begin{longtable}{@{}p{0.30\textwidth}p{0.52\textwidth}@{}}
\toprule
ID & Resultado \\
\midrule
\endhead
TEST-INT-NOTIF-M13-012 & Marca todas as ativas do UID, inclusive outro papel, sem \texttt{DELETE} \\
TEST-INT-NOTIF-M13-013 & Retry preserva \texttt{lida\_em} já gravado e o de item já lido \\
TEST-INT-NOTIF-M13-014 & Corte estável exclui emissão posterior; nova invocação a marca \\
TEST-INT-NOTIF-M13-015 & Não age na caixa de terceiro \\
TEST-INT-NOTIF-M13-016 & Paginação reentrante com o mesmo corte (2 + 1) \\
TEST-INT-NOTIF-M13-017 & Documento com \texttt{id\_destinatario} divergente não é marcado \\
TEST-INT-NOTIF-M13-018 & \texttt{lida=false} com \texttt{lida\_em} preenchido falha fechado \\
TEST-INT-NOTIF-M13-019 & Autoridade inativa nega sem marcar nada \\
TEST-INT-NOTIF-M13-020 & Corte inválido é \texttt{invalid-argument} \\
TEST-RULES-NOTIF-010 & Listagem filtrada por \texttt{id\_destinatario} é permitida \\
TEST-RULES-NOTIF-011 & Listagem sem o filtro de coerência é negada \\
\bottomrule
\end{longtable}
\normalsize

RED: a suíte de integração não compilava contra a implementação anterior
(\texttt{limparTudoNotificacoes} inexistente), caracterizando a feature ausente.
GREEN sob Node 20.19.6: notificações 31/31 (integração 20/20 e Rules 11/11);
unit 71/71; Rules 77/77. Regressão: \texttt{build} PASS e integração completa
138 passam / 14 falham (falhas legadas em reagentes/patrimônio/roteiros,
inalteradas).

\textbf{Estado final:} \texttt{IMP-NOTIF-003} = \texttt{NÃO DIVERGENTE} no
escopo backend auditado; UI de disparo permanece sob \texttt{IMP-NOTIF-001} e
\texttt{IMP-UI-004}.

\subsection{Hardening do conjunto ativo (RN-M13-03)}

Auditoria independente reabriu \texttt{IMP-NOTIF-003}: ``Limpar tudo'' marcava
qualquer notificação \texttt{lida=false} do corte, sem considerar
\texttt{expira\_em}. Uma notificação já expirada não pertence ao conjunto ativo
(RN-M13-03; Alloy \texttt{ativo = notExpira != EVencido},
\texttt{ExpiradoForaDoAtivo}) e não pode ser tratada como elegível.

\paragraph{Correção.} A consulta continua com corte estável
(\texttt{emitida\_em <= corte}) e passa a ignorar, na rodada, a notificação cujo
\texttt{expira\_em} já foi alcançado (avaliado no instante de processamento);
o fato permanece persistido. Como itens ignorados não são mutados, o avanço usa
o cursor estável \texttt{(emitida\_em, docId)} — que também preserva empates do
\texttt{serverTimestamp} e evita re-scan infinito —, devolvido em
\texttt{proximo\_cursor}. \texttt{expira\_em=null} (ex.: \texttt{ESCASSEZ\_ESTOQUE})
permanece sempre elegível. Nenhum índice novo é necessário: o índice
\texttt{Notificacoes(id\_destinatario, lida, emitida\_em)} já cobre a consulta.

\small
\begin{longtable}{@{}p{0.30\textwidth}p{0.52\textwidth}@{}}
\toprule
ID & Resultado \\
\midrule
\endhead
TEST-INT-NOTIF-M13-021 & Notificação expirada não é marcada e permanece persistida \\
TEST-INT-NOTIF-M13-022 & Expiração futura é elegível \\
TEST-INT-NOTIF-M13-023 & \texttt{expira\_em=null} permanece elegível \\
TEST-INT-NOTIF-M13-024 & Paginação com expiradas intercaladas não perde ativos (cursor) \\
TEST-INT-NOTIF-M13-025 & Retry com mesmo cursor preserva \texttt{lida\_em} e não duplica \\
TEST-INT-NOTIF-M13-026 & Cursor inválido é \texttt{invalid-argument} \\
\bottomrule
\end{longtable}
\normalsize

RED: os casos 021, 024, 025 e 026 falharam contra a implementação anterior
(marcava expirada, não tinha cursor). GREEN sob Node 20.19.6: notificações
37/37 (integração 26/26 e Rules 11/11); unit 71/71; Rules 77/77; integração
completa 151 passam / 14 falham (legadas em reagentes/patrimônio/roteiros).

\textbf{Estado final:} \texttt{IMP-NOTIF-003} = \texttt{NÃO DIVERGENTE} no
escopo backend auditado.
